\begin{abstract}
Coding tasks reconstructed from agent sessions are increasingly used as evaluation and training environments. In a real session the user's requirements evolve: later turns correct, refine or retract the initial request. If a task's executable verifier is built from an incomplete account of those requirements, it can reward solutions the user would have rejected (false acceptance) or reject solutions the user would have accepted (false rejection). We study this \emph{reward-fidelity} failure directly. We compare verifiers built from the initial request with verifiers built from an evidence-linked, correction-aware requirement contract; each is evaluated with and without a predefined audit procedure, against repository tests and a full-session LLM baseline. All verifiers score the \emph{same} held-out candidate patches, including the agent's naturally observed pre-correction code state, constructed near misses, and independent valid alternatives. Grading runs in a process-separated sandbox whose isolation properties are tested. \unfilled{summary of results; the current draft reports only an automated pilot on controlled sessions with model-proposed labels}.
\end{abstract}
